%% ECON 282E -- Homework 2, Part A: Heterogeneity and Learning (core, S6-S7)
%% Component 2 of the project-manager ladder (HW1 -> HW2 -> final project).
%% Part B (hw2b.tex) adds S8-S9, the student's route, and the peer audit.
%% Build:  python3 tools/build_docs.py homeworks-exams/HW2/hw2a.tex
%% With answers:  python3 tools/build_docs.py --solutions homeworks-exams/HW2/hw2a.tex
%% Answers live in solutions/ (gitignored); solution PDFs are gitignored too.
%% Neither must ever be committed.

\documentclass[11pt]{article}

\newif\ifsolutions \solutionsfalse

\usepackage[margin=1in]{geometry}
\usepackage{amsmath,amssymb,booktabs,enumitem,xcolor}
\usepackage[T1]{fontenc}
\usepackage{microtype}
\usepackage[colorlinks=true,linkcolor=blue!50!black,urlcolor=blue!50!black]{hyperref}

\newcommand{\solution}[1]{\ifsolutions\medskip
  {\color{blue!45!black}\textbf{Solution.} #1}\medskip\fi}

%% The answer key lives in solutions/ (gitignored, never published), so this
%% source carries no answers. --solutions pulls it in; a plain build ignores it.
\AtBeginDocument{\ifsolutions
  \IfFileExists{solutions/hw2a_key.tex}{\input{solutions/hw2a_key.tex}}
    {\GenericWarning{}{Answer key solutions/hw2a_key.tex not found}}\fi}

\title{\vspace{-2em}\textbf{ECON 282E --- Homework 2, Part A}\\[2mm]
       \large Heterogeneity and Learning: the Distribution, and a Network You Can Trust}
\author{Foundations of Macroeconomics \\ University of California, Riverside}
\date{Out: Thursday, October 29, 2026 \quad\textbullet\quad Due (with Part B): Monday, November 23, 2026}

\begin{document}
\maketitle
\thispagestyle{empty}

\noindent\rule{\linewidth}{0.4pt}\par
\small
\textbf{Where this sits.} Homework 1 built \textbf{component 1}: a calibrated representative-agent
engine, a data pipeline, and a library of solvers you have tested. Homework 2 builds
\textbf{component 2}: a heterogeneous-agent economy and one learning or measurement module. The final
project \textbf{assembles} them to answer your team's question. So reuse your own HW1 code wherever
you can --- and when you do, say what you reused and how you re-checked it.

\textbf{More of the lead is yours.} In HW1 the questions were set for you. Here a \textbf{common core}
(Part A, and C3--C4 in Part B) makes sure everyone can use every method of Sessions 6--9. Then you
choose a \textbf{route} (Part B) and pose your own question inside it. Keep your HW1 charter going:
update it for this homework, and hand in the new version.

\textbf{Calibration.} Use \textbf{your} HW1 values of $\alpha$, $\beta$ and $\delta$ (quarterly). If
you replaced them in HW1, use the replacement and say so. Add CRRA utility with $\sigma=2$, a borrowing
limit $a'\ge0$, and an earnings process $\ln z'=\rho_z\ln z+\varepsilon'$ with $\rho_z=0.95$ and an
unconditional standard deviation of $\ln z$ of $0.30$, discretized by Rouwenhorst with 7 states. You
may argue for another process --- then show the numbers both ways.

\textbf{What to hand in (with Part B).} One PDF; a repository whose \texttt{run\_all.py} reproduces
every number and figure; the updated charter; and \texttt{AI\_LOG.md} with its validation ledger.
The HW1 templates still apply (\texttt{homeworks-exams/HW1/templates/}).
\normalsize\par
\noindent\rule{\linewidth}{0.4pt}\par

%=============================================================================
\section*{C1. The Aiyagari economy, end to end (25 points)}
\emph{Session 6, Block A.}

\begin{enumerate}[label=(\alph*), leftmargin=*]

\item \textbf{(5 pts) The equilibrium, on paper.} Define the stationary recursive competitive
equilibrium (you did this in words in HW1 A1(d); now write it formally, with the Markov operator $Q$
and the invariant measure $\lambda$). Say which objects are functions, which are numbers, and which
is a measure.

\solution{\keyCone}

\item \textbf{(10 pts) Solve it as a nested fixed point.} Inner loop: the household's policy at a
given $r$ (EGM, or your best HW1 solver). Middle loop: the invariant distribution. Outer loop: the
interest rate that clears the capital market. Report $r^*$ (quarterly and annual), $K$, $K/Y$, the
wealth Gini, the top-10\% wealth share, the fraction at the borrowing limit, and the number of
iterations in each loop. Plot capital supply and demand against $r$, with $1/\beta-1$ marked.

\solution{\keyCtwo}

\item \textbf{(10 pts) Quality control.} Before anyone else trusts these numbers, show that:
(i) the market clears to a stated tolerance; (ii) $\lambda$ sums to one and puts negligible mass at
the top of the asset grid; (iii) the headline numbers are stable when you double the number of asset
grid points --- tabulate $N_a\in\{200,400,800\}$. Which number is the \emph{least} stable, and why?
Finally, compare $r^*$ and $K/Y$ with your HW1 representative-agent economy at the same
$\alpha,\beta,\delta$, and relate the difference to Session~2.

\solution{\keyCthree}

\end{enumerate}

%=============================================================================
\section*{C2. Learn it, then prove the network right (20 points)}
\emph{Session 7.}

Train a PyTorch network on the Euler residual (7.B1) for \textbf{one} of two targets --- your choice,
stated in the charter:
\begin{itemize}[leftmargin=*,itemsep=0pt]
  \item[(A)] \textbf{Your HW1 stochastic growth model.} Your HW1 Chebyshev collocation solution is the
        benchmark.
  \item[(B)] \textbf{The Aiyagari household at your $r^*$ from C1.} Your C1 EGM policy is the
        benchmark. The borrowing limit makes this harder: handle it with a Fischer--Burmeister penalty
        (2.A1, 4.A2) or by building the constraint into the architecture (7.B3), and say which.
\end{itemize}

\begin{enumerate}[label=(\alph*), leftmargin=*]

\item \textbf{(8 pts) Train to a plateau, not to a budget.} Report the architecture, the parameter
count, how you enforce feasibility, how the expectation is computed, the optimizer and its schedule,
and the loss curve. How do you know training has stopped improving?

\solution{\keyCfour}

\item \textbf{(8 pts) Grade it against the benchmark.} Unit-free Euler residuals at 2{,}000 off-grid
points over the whole domain and over its interior 90\%, and the maximum policy gap to the benchmark.
\emph{Where} is the network's worst residual --- at the boundary of the training region, near the
constraint, or spread out? Plot the residual against the state.

\solution{\keyCfive}

\item \textbf{(4 pts) The decision.} For this model, is the network worth it? Name one model in this
course where you would \emph{not} use the classical method you just beat or lost to, and say why.

\solution{\keyCsix}

\end{enumerate}

%=============================================================================
\section*{C0. Red-team your AI assistant (5 points)}

Ask an AI assistant for \textbf{one} of: the SRCE definition of C1(a); the Young (2010) lottery that
moves mass between grid points; or the loss function of C2. Find and document a real error, an
unstated assumption, or a silent change of notation in what it gives you, and show the check that
caught it (a derivation, a unit test, a counterexample). Put the exchange and the check in
\texttt{AI\_LOG.md}. \emph{If you find nothing wrong, show the test you ran that would have caught
an error --- and then try a harder request.}

\solution{\keyCseven}

\vfill
\noindent\rule{\linewidth}{0.4pt}\par
\small
\textbf{Reporting standard.} With every computed number: the method, the grid, the tolerance of
\emph{each} loop, the dtype, the hardware and the wall time; with every network: the seed, the
architecture, the training budget and the validation points. A number without these is not yet a
result. \textbf{Part A: 50 points.}
\normalsize

\end{document}
